\documentclass{article}
\usepackage{amsmath,multicol}
\usepackage{amsthm}
\usepackage{amssymb}
\usepackage{amsfonts}
\usepackage{subcaption}
\usepackage{graphicx}
\usepackage{latexsym}
\usepackage{enumerate}
\usepackage{cite, hyperref, xcolor}
\usepackage{bm}
\usepackage{tikzsymbols}
\usepackage{multicol}
\usepackage{tasks}
\usepackage{mathdots}
\usepackage[mathscr]{eucal}
\usepackage{comment}
\usepackage{tikz-cd}
%\usepackage{extsizes}
%\usepackage{graphicx}

\settasks{counter-format = (tsk[a]), label-offset = {.7em} }

\usepackage[sc]{mathpazo}
%\usepackage{kpfonts}
\linespread{1.05}
%\usepackage[T1]{fontenc}
\usepackage{tikz}

%%  Renewed
\renewcommand{\Re}[1]{\operatorname{Re} #1 }
\renewcommand{\Im}[1]{\operatorname{Im} #1}

%% New Commands
\newcommand{\inv}{^{-1}}
\newcommand{\X}{\mathscr{X}} % fcnl analysis, vector space X
\newcommand{\C}{\mathbb{C}} % complex numbers
\newcommand{\A}{\mathcal{A}} % sigma-algebra
\newcommand{\B}{\mathcal{B}} % alternate sigma-algebra
\newcommand{\h}{\mathcal{H}} % hilbert series
\newcommand{\T}{\mathscr{T}} % topology
\renewcommand{\S}{\mathscr{S}} % alternate topology
\newcommand{\<}{\langle} % gen set left
\renewcommand{\>}{\rangle} % gen set right
\newcommand{\Z}{\mathbb{Z}} % integers
\newcommand{\N}{\mathbb{N}} % naturals
\newcommand{\R}{\mathbb{R}} % reals
\newcommand{\LL}{\mathcal{L}} % Lebesgue measurable sets
\newcommand{\RR}{\mathcal{R}} % ring?
\newcommand{\D}{\mathbb{D}} % open unit disk
\newcommand{\Q}{\mathbb{Q}} % rationals
\newcommand{\F}{\mathscr{F}} % family of functions
\newcommand{\cl}{\operatorname{cl}}
\newcommand{\ran}{\operatorname{ran}}
\newcommand{\norm}[1]{\| #1 \|} % norm
\newcommand{\inner}[1]{\langle #1 \rangle} %inner product or gen set
\newcommand{\fp}{\mathfrak p} %prime ideal p notation
\newcommand{\fm}{\mathfrak m} %max ideal m notation
\providecommand{\ceil}[1]{\left \lceil #1 \right \rceil}
\providecommand{\floor}[1]{\left \lfloor #1 \right \rfloor}
\renewcommand{\vec}[1]{{\bm #1}}
\newcommand{\eps}{\varepsilon} % comment out if important people want to use this preamble hehe

\DeclareMathOperator{\ord}{ord}
\DeclareMathOperator{\Gal}{Gal}
\DeclareMathOperator{\Hom}{Hom}
\DeclareMathOperator{\End}{End}
\DeclareMathOperator{\Aut}{Aut}
\DeclareMathOperator{\GL}{GL}
\DeclareMathOperator{\Ann}{Ann}
\DeclareMathOperator{\Tor}{Tor}

%%  Matrices
\newcommand{\minimatrix}[4]{\begin{bmatrix} #1 & #2 \\ #3 & #4 \end{bmatrix}  }
\newcommand{\megamatrix}[9]{\begin{bmatrix} #1 & #2 & #3 \\ #4 & #5 & #6 \\ #7 & #8 & #9\end{bmatrix}  }
\newcommand{\tworowvector}[2]{[#1\,\,#2]}

%\renewcommand{\labelenumi}{(\roman{enumi})}

\newcommand{\diag}{\operatorname{diag}}
\newcommand{\rank}{\operatorname{rank}}
\newcommand{\tr}{\operatorname{tr}}
\newcommand{\twovector}[2]{\begin{bmatrix} #1\\#2 \end{bmatrix} }
\newcommand{\fourrowvector}[4]{\begin{bmatrix} #1\,\, #2\,\, #3\,\, #4 \end{bmatrix} }
%\newcommand{\threevector}[3]{\begin{bmatrix} #1\\#2\\#3 \end{bmatrix} }

%\newcommand{\comment}[1]{\marginpar{\tiny\color{blue}$\bullet$  #1}}
%\newcommand{\scomment}[1]{\marginpar{\tiny\color{red}$\bullet$  #1}}

%%%
%%% Theorem Styles
%%%
% \numberwithin{equation}{section}
\theoremstyle{plain}
\newtheorem{Theorem}[equation]{Theorem}
\newtheorem{Conjecture}[equation]{Conjecture}
\newtheorem{Lemma}[equation]{Lemma}
\newtheorem{Proposition}[equation]{Proposition}
\newtheorem{Corollary}[equation]{Corollary}
\theoremstyle{remark}
\newtheorem{Remark}[equation]{Remark}
\theoremstyle{definition}
\newtheorem{Definition}[equation]{Definition}
\newtheorem{Example}[equation]{Example}
\newcounter{question}
\newtheorem{Question}[question]{Question}
\newtheorem{examproblem}[equation]{Problem}

\newcommand{\spatula}{
\vspace*{3cm}
\begin{tikzpicture}[scale = 0.3]
\draw[line width=0.6mm,-] (0,-1)--(0,-0.5);
\draw[-] (0,-0.5) -- (0,0) -- (0.3536,0.3536) -- (0.3536,0.8536) -- (-0.3536,0.8536) -- (-0.3536,0.3536) -- (0,0) -- cycle;
\draw[-] [line width = 0.25mm,-] (-0.1768,0.3536)--(-0.1768,0.6869);
\draw[-] [line width = 0.25mm,-] (0,0.3536)--(0,0.6869);
\draw[-] [line width = 0.25mm,-] (0.1768,0.3536)--(0.1768,0.6869);
\end{tikzpicture}
}
\newcommand{\rockon}{
\vspace*{3cm}
\begin{tikzpicture}[scale=0.075]
\draw (0,-.25) -- (0,2.25) arc (180:0:.5) -- (1,.75) % index
arc (180:0:0.5) % middle
arc (180:0:.5) % ring
-- (3,1.75) arc (180:0:.5) -- (4,-1.25) % pinkie
arc (360:270:1.75) -- (1.5,-3) arc (270:180:1.75) % palm
--(-.25,-1) arc (180:90:1) -- (1.5,0) arc (450:270:.5) -- (1,-1) arc (90:270:.15) % thumb
arc(90:0:1); % inside palm
\draw (1,.75) -- (1,0); % middle 2
\draw (2,.75) -- (2,-.5) arc (180:360:.5) -- (3,.75); %ring 2
\draw (.5,-2.7)--(.5,-3.3); % lwrist
\draw (3.25,-2.7)--(3.25,-3.3); % rwrist
\draw (0,-4.3) rectangle (3.75,-3.3); %sleeve
\end{tikzpicture}
}

\renewcommand\qedsymbol{$_{\spatula}$}
%\renewcommand\qedsymbol{$_{\rockon}$}

%%%% HOMEWORK
\newenvironment{exercise}[2][Exercise]{\begin{trivlist}
\item[\hskip \labelsep {\bfseries #1}\hskip \labelsep {\bfseries #2.}]}{\end{trivlist}}

\allowdisplaybreaks
%%
%% MAIN DOCUMENT
%%

\textwidth=6.5in
\textheight=8.5in
\topmargin=-1in
\headheight=0.5in
\headsep=.5in
\hoffset= -1in


%%%%%%% numbering stuff
%\setcounter{secnumdepth}{0}
\usepackage{chngcntr}

\counterwithin*{equation}{section}
\counterwithin*{equation}{subsection}


\begin{document}

\title{A lot of algebra prelims}
\date{}
\author{}


\maketitle
%Paste homework .tex here
\tableofcontents

\pagebreak

\section{2018}
\begin{examproblem}[2018, 1]
	Show that no finite group is the union of conjugates of a proper subgroup.
\end{examproblem}

\begin{examproblem}[2018, 2]
	Classify all groups of order 18 up to isomorphism.
\end{examproblem}

\begin{examproblem}[2018, 3]
	Let $\alpha,\beta$ denote the unique positive
	real $5^{\text{th}}$ root of 7 and $4^{\text{th}}$ root of 5,
	respectively. Determine the degree of $\mathbb Q(\alpha,\beta)$ over $\mathbb Q$.
\end{examproblem}

\begin{examproblem}[2018, 4]
	Show that the field extension $\mathbb Q\subseteq\mathbb Q\left(
	\sqrt{2+\sqrt2}\right)$ is Galois and determine its Galois group.
\end{examproblem}

\begin{examproblem}[2018, 5]
	Let $M$ be a square matrix over a field $K$. Use a suitable canonical
	form to show that $M$ is similar to its transpose $M^T$.
\end{examproblem}



\begin{examproblem}[2018, 7]
	Let $G$ be a finite group and $\pi$, $\pi'$ be two irreducible
	representations of $G$. Prove or disprove the following assertion:
	$\pi$ and $\pi'$ are equivalent if and only if $\det\pi(g)=\det\pi'(g)$
	for all $g\in G$.
\end{examproblem}

\pagebreak
%%%%%%%%%%%
\section{2017}
\begin{examproblem}[2017, 1]
	Let $R$ be a Noetherian ring. Prove that $R[x]$ and $R[[x]]$ are both
	Noetherian. (The first part of the question is asking you to prove the
	Hilbert Basis Theorem, not to use it!)
\end{examproblem}

\begin{examproblem}[2017, 2]
	Classify (with proof) all fields with finitely many elements.
\end{examproblem}

\begin{examproblem}[2017, 3]
	Suppose $A$ is a commutative ring and $M$ is a finitely presented module.
	Given any surjection $\phi:A^n\rightarrow M$ from a finite free
	$A$-module, show that $\ker\phi$ is finitely generated.
\end{examproblem}

\begin{examproblem}[2017, 4.]
	Classify all groups of order 57.
\end{examproblem}

\begin{examproblem}[2017, 7]
	Show that a finite simple group cannot have a 2-dimensional irreducible
	representation over $\mathbb C$. (Hint: the determinant might prove useful.)
\end{examproblem}

\pagebreak
%%%%%%%%%%%

\section{2016}
\begin{examproblem}[2016, 1.]
	Let $G$ be a finite simple group. Assume that every proper subgroup
	of $G$ is abelian. Prove that then $G$ is cyclic of prime order.
\end{examproblem}


\begin{examproblem}[2016, 2]
	Let $a\in\mathbb N$, $a>0$. Compute the Galois group of the splitting field
	of the polynomial $x^5-5a^4x+a$ over $\mathbb Q$.
\end{examproblem}

\begin{examproblem}[2016, 4.] 
  Recall that an inner automorphism of a group is an automorphism given by conjugation by an element of the group. An outer automorphism is an automorphism that is not inner.
	\begin{enumerate}[(a)]
		\item Prove that $S_5$ has a subgroup of order 20.
		\item Use the subgroup from (a) to construct a degree 6 permutation representation of $S_5$ (i.e., an embedding $S_5 \hookrightarrow S_6$ as a transitive permutation group on 6 letters).
		\item Conclude that $S_6$ has an outer automorphism.
	\end{enumerate}
\end{examproblem}

\begin{examproblem}[2016, 5.]
    Let $A$ be a commutative ring and $M$ a finitely generated $A$-module. Define
    \[ \Ann(M) = \{a \in A: am = 0 \text{ for all } m \in M\}.\]
    Show that for a prime ideal $\mathfrak p \subset A$, the following are equivalent:
    \begin{enumerate}[(a)]
        \item $\Ann(M) \not\subset \mathfrak p$
        \item The localization of $M$ at the prime ideal $\mathfrak p$ is $0$.
        \item $M \otimes_A k(\mathfrak p) = 0$, where $k(\mathfrak p) = A_{\mathfrak p}/\mathfrak p A_{\mathfrak p}$ is the residue field of $A$ at $\mathfrak p$.
    \end{enumerate}
\end{examproblem}

\begin{examproblem}[2016, 6.]
  Let $A = \C[x,y]/(y^2-(x-1)^3 - (x-1)^2)$.
  \begin{enumerate}[(a)]
      \item Show that $A$ is an integral domain and sketch the $\R$-points of $\text{Spec} A$.
      \item Find the integral closure of $A$. Recall that for an integral domain $A$ with fraction field $K$, the integral closure of $A$ in $K$ is the set of all elements of $K$ integral over $A$.
  \end{enumerate}
\end{examproblem}

\begin{examproblem}
    Let $R = k[x,y]$ where $k$ is a field, and let $I=(x,y)R$.
    \begin{enumerate}
        \item Show that
        \[
        \begin{tikzcd}
        0 \ar[r] & R \arrow[r,"\phi"] & R \oplus R \arrow[r,"\psi"] & R \arrow[r] & k \arrow[r] & 0
        \end{tikzcd}
        \]
        where $\phi(a) = (-ya,xa)$, $\psi((a,b)) = xa+yb$ for $a,b \in R$, is a projective resolution of the $R$-module $k \simeq R/I$.
        \item Show that $I$ is not a flat $R$-module by computing $\Tor_i^R(I,k)$
    \end{enumerate}
\end{examproblem}


\pagebreak
%%%%%%%%%%%

\section{2015}

\begin{examproblem}[2015, 1]
	\begin{enumerate}[(a)]
		\item Find an irreducible polynomial of degree 5 over the
			field $\mathbb Z/2$ of two elements and use it to construct
			a field of order 32 as a quotient of the
			polynomial ring $\mathbb Z/2[x]$.
		\item Using the polynomial found in part (a), find a $5\times5$
			matrix $M$ over $\mathbb Z/2$ of order 31, so that $M^{31}=I$
			but $M\neq I$.
	\end{enumerate}
\end{examproblem}

\begin{examproblem}[2015, 2]
	Find the minimal polynomial of $\sqrt2+\sqrt3$ over $\mathbb Q$.
	Justify your answer.
\end{examproblem}



\begin{examproblem}[2015, 3]
	\begin{enumerate}[(a)]
		\item Let $R$ be a commutative ring with no nonzero nilpotent
			elements. Show that the only units in the polynomial ring
			$R[x]$ are the units of $R$, regarded as constant polynomials.
		\item Find all units in the polynomial ring $\mathbb Z_4[x]$.
	\end{enumerate}
\end{examproblem}

\begin{examproblem}[2015, 4.]
	Let $p$, $q$ be two distinct primes. Prove that there is at most one
	non-abelian group of order $pq$ and describe the pairs $(p,q)$
	such that there is no non-abelian group of order $pq$.
\end{examproblem}


\begin{examproblem}[2015, 5]
	\begin{enumerate}[(a)]
		\item Let $L$ be a Galois extension of a field $K$ of degree 4.
			What is the minimum number of subfields there could
			be strictly between $K$ and $L$? What is the maximum
			number of such subfields? Give examples where these
			bounds are attained.
		\item How do these numbers change if we assume only that $L$
			is separable (but not necessarily Galois) over $K$?
	\end{enumerate}

\end{examproblem}

\begin{examproblem}[2015, 6]
	Let $R$ be a commutative algebra over $\mathbb C$. A derivation of $R$ is a
	$\mathbb C$-linear map $D:R\rightarrow R$ such that (i) $D(1)=0$ and
	(ii) $D(ab)=D(a)b+aD(b)$ for all $a,b\in R$.
	\begin{enumerate}[(a)]
		\item Describe all derivations of the polynomial ring $\mathbb C[x]$.
		\item Let $A$ be the subring (or $\mathbb C$-subalgebra) of
			$\End_{\mathbb C}(\mathbb C[x])$ generated by all derivations
			of $\mathbb C[x]$ and the left multiplications by $x$.
			Prove that $\mathbb C[x]$ is a simple left $A$-module.
			Note that the inclusion $A\rightarrow\End_{\mathbb C}(\mathbb C[x])$
			defines a natural left $A$-module structure on $\mathbb C[x]$.
	\end{enumerate}
\end{examproblem}


\begin{examproblem}[2015, 7]
	Let $G$ be a non-abelian group of order $p^3$ with $p$ a prime.
	\begin{enumerate}[(a)]
		\item Determine the order of the center $Z$ of $G$.
		\item Determine the number of inequivalent complex
			1-dimensional representations of $G$.
		\item Compute the dimensions of all the inequivalent irreducible
			representations of $G$ and verify that the number of such
			representations equals the number of conjugacy classes of
			$G$.
	\end{enumerate}
\end{examproblem}

\newpage

\section{2014}

\begin{examproblem}[2014, 1.]
	\begin{enumerate}[(a)]
		\item Let $G$ be a group (not necessarily finite) that contains
			a subgroup of index $n$. Show that $G$ contains a
			\textit{normal} subgroup $N$ such that $n\leq[G:N]\leq n!$
		\item Use part (a) to show that there is no simple group
			of order 36.
	\end{enumerate}
\end{examproblem}


\begin{examproblem}[2014, 2]
	Let $p$ be a prime, let $\mathbb F_p$ be the $p$-element field,
	and let $K=\mathbb F_p(t)$ be the field of rational functions in
	$t$ with coefficients in $\mathbb F_p$. Consider the polynomial $f(x)=
	x^p-t\in K[x]$.
	\begin{enumerate}[(a)]
		\item Show that $f$ does not have a root in $K$.
		\item Let $E$ be the splitting field of $f$ over $K$.
			Find the factorization of $f$ over $E$.
		\item Conclude that $f$ is irreducible over $K$.
	\end{enumerate}
\end{examproblem}

\begin{examproblem}[2014, 3]
	Recall that a ring $A$ is called \textit{graded} if it admits a direct
	sum decomposition $A=\oplus_{n=0}^{\infty}A_n$ as abelian groups, with the
	property that $A_iA_j\subseteq A_{i+j}$ for all $i,j\geq 0$.
	Prove that a graded commutative ring $A=\oplus_{n=0}^{\infty} A_n$
	is Noetherian if and only if $A_0$ is Noetherian and $A$ is finitely
	generated as an algebra over $A_0$.
\end{examproblem}

\begin{examproblem}[2014, 4]
	Let $R$ be a ring with the property that $a^2=a$ for all $a\in R$.
	\begin{enumerate}[(a)]
		\item Compute the Jacobson radical of $R$.
		\item What is the characteristic of $R$?
		\item Prove that $R$ is commutative.
		\item Prove that if $R$ is finite, then $R$ is isomorphic
			(as a ring) to $(\mathbb Z/2\mathbb Z)^d$ for some $d$.
	\end{enumerate}
\end{examproblem}

\begin{examproblem}[2014, 6]
	Let $\overline{\mathbb F_p}$ denote the algebraic closure of $\mathbb F_p$. Show that
	the Galois group $\Gal(\overline{\mathbb F_p}/\mathbb F_p)$
	has no non-trivial finite subgroups.
\end{examproblem}

\begin{examproblem}[2014, 7]
	Let $C_p$ denote the cyclic group of order $p$.
	\begin{enumerate}[(a)]
		\item Show that $C_p$ has two irreducible representations over
			$\mathbb Q$ (up to isomorphism), one of dimension 1
			and one of dimension $p-1$.
		\item Let $G$ be a finite group, and let $\rho:G\rightarrow
			\GL_n(\mathbb Q)$ be a representation of $G$ over $\mathbb Q$.
			Let $\rho_{\mathbb C}:G\rightarrow\GL_n(\mathbb C)$ denote
			$\rho$ followed by the inclusion $\GL_n(\mathbb Q)\rightarrow
			\GL_n(\mathbb C)$. Thus $\rho_{\mathbb C}$ is a representation
			of $G$ over $\mathbb C$, called the \textit{complexification}
			of $\rho$. We say that an irreducible representation $\rho$
			of $G$ is \textit{absolutely irreducible} if its
			complexification remains irreducible over $\mathbb C$.\\
			Now suppose $G$ is abelian and that every representation
			of $G$ over $\mathbb Q$ is absolutely irreducible. Show that
			$G\cong(C_2)^k$ for some $k$ (i.e., is a product of
			cyclic groups of order 2).
	\end{enumerate}
\end{examproblem}

\begin{examproblem}[2014, 8]
	Let $G$ be a finite group and $\mathbb Z[G]$ the internal group algebra.
	Let $\mathcal Z$ be the center of $\mathbb Z[G]$. For each conjugacy class
	$C\subseteq G$, let $P_C=\sum_{g\in C}g$.
	\begin{enumerate}[(a)]
		\item Show that the elements $P_C$ form a $\mathbb Z$-basis for
			$\mathcal Z$. Hence $\mathcal Z\cong\mathbb Z^d$ as an
			abelian group, where $d$ is the number of
			conjugacy classes in $G$.
		\item Show that if a ring $R$ is isomorphic to $\mathbb Z^d$ as
			an abelian group, then every element in $R$ satisfies
			a monic integral polynomial.
			(\textbf{Hint:} Let $\{v_1,\dots,v_d\}$ be a basis of
			$R$ and for a fixed non-zero $r\in R$, write $rv_i=\sum_j
			a_{ij}v_j$. Use the Hamilton-Cayley theorem.)
		\item Let $\pi:G\rightarrow\GL(V)$ be an irreducible representation
			of $G$ (over $\mathbb C$). Show that $\pi(P_C)$ acts on $V$
			as multiplication by the scalar
			$$
			\frac{|C|\chi_{\pi}(C)}{\dim V},
			$$
			where $\chi_{\pi}(C)$ is the value of the character
			$\chi_{\pi}$ on any element of $C$.
		\item Conclude that $|C|\chi_{\pi}(C)/\dim V$ is an algebraic
			integer.
	\end{enumerate}
\end{examproblem}

\newpage

\section{2013}

\begin{examproblem}[2013, 3.]
	\begin{enumerate}[(a)]
		\item Suppose that $G$ is a finitely generated group. Let $n$
			be a positive integer. Prove that $G$ has only finitely
			many subgroups of index $n$
		\item Let $p$ be a prime number. If $G$ is any finitely-generated
			abelian group, let $t_p(G)$ denote the number of
			subgroups of $G$ of index $p$. Determine the possible
			values of $t_p(G)$ as $G$ varies over all
			finitely-generated abelian groups.
	\end{enumerate}
\end{examproblem}

\begin{examproblem}[2013, 4.]
	Suppose that $G$ is a finite group of order 2013. Prove that $G$ has
	a normal subgroup $N$ of index 3 and that $N$ is a cyclic group.
	Furthermore, prove that the center of $G$ has order divisible by 11.
	(You will need the factorization $2013=3\cdot11\cdot61$.)
\end{examproblem}

\begin{examproblem}[2013, 6]
	This question concerns an extension $K$ of $\mathbb Q$ such that
	$[K:\mathbb Q]=8$. Assume that $K/\mathbb Q$ is Galois and let
	$G=\Gal(K/\mathbb Q)$. Furthermore, assume that $G$ is non-abelian.
	\begin{enumerate}[(a)]
		\item Prove that $K$ has a unique subfield $F$ such that
			$F/\mathbb Q$ is Galois and $[F:\mathbb Q]=4$.
		\item Prove that $F$ has the form $F=\mathbb Q(\sqrt{d_1},\sqrt{d_2})$
			where $d_1,d_2$ are non-zero integers.
		\item Suppose that $G$ is the quaternionic group. Prove
			that $d_1$ and $d_2$ are positive integers.
	\end{enumerate}
\end{examproblem}

\begin{examproblem}[2013, 8]
	This question concerns the polynomial ring $R=\mathbb Z[x,y]$ and the ideal
	$I=(5,x^2+2)$ in $R$.
	\begin{enumerate}[(a)]
		\item Prove that $I$ is a prime ideal of $R$ and that $R/I$ is
			a PID.
		\item Give an explicit example of a maximal ideal of $R$ which
			contains $I$. (Give a set of generators for such an ideal.)
		\item Show that there are infinitely many distinct maximal ideals
			in $R$ which contain $I$.
	\end{enumerate}
\end{examproblem}

\newpage

\section{2012}

\begin{examproblem}[2012, 1.]
	Classify all groups of order 2012 up to isomorphism. (Hint: 503 is prime).
\end{examproblem}

\begin{examproblem}[2012, 2.]
	For any positive integer $n$, let $G_n$ be the group generated by
	$a$ and $b$ subject to the following three relations:
	$$
	a^2=1, \quad b^2=1, \quad \text{and} \quad (ab)^n=1.
	$$
	\begin{enumerate}[(a)]
		\item Find the order of the group $G_n$
	\end{enumerate}
\end{examproblem}
(We don't know how to do the rest of the problem)

\begin{examproblem}[2012, 6]
	Determine the Galois groups of the following polynomials over $\mathbb Q$.
	\begin{enumerate}[(a)]
		\item $f(x)=x^4+4x^2+1$
		\item $f(x)=x^4+4x^2-5$.
	\end{enumerate}
\end{examproblem}

\begin{examproblem}[2012, 3]
	Let $R$ be a (commutative) principal ideal domain, let $M$ and $N$ be
	finitely generated free $R$-modules, and let $\varphi:M\rightarrow N$ be
	an $R$-module homomorphism.
	\begin{enumerate}[(a)]
		\item Let $K$ be the kernel of $\varphi$. Prove that
			$K$ is a direct summand of $M$.
		\item Let $C$ be the image of $\varphi$. Show by example
			(specifying $R$, $M$, $N$, and $\varphi$) that $C$
			need not be a direct summand of $N$.
	\end{enumerate}
\end{examproblem}

\newpage

\section{2011}

\begin{examproblem}[2011, 2.]
	In this problem, as you apply Sylow's Theorem, state precisely which portions
	you are using.
	\begin{enumerate}[(a)]
		\item Prove that there is no simple group of order 30.
		\item Suppose that $G$ is a simple group of order 60. Determine
			the number of $p$-Sylow subgroups of $G$ for each prime
			$p$ dividing 60, then prove that $G$ is isomorphic to
			the alternating group $A_5$.
	\end{enumerate}
	Note: in the second part, you needn't show that $A_5$ is simple.
	You need only show that if there is a simple group of order 60, then
	it must be isomorphic to $A_5$.
\end{examproblem}

\begin{examproblem}[2011, 3]
	Describe the Galois group and the intermediate fields of the cyclotomic
	extension $\mathbb Q(\zeta_{12})/\mathbb Q$.
\end{examproblem}



\begin{examproblem}[2011, 4]
	Let
	$$
	R=\mathbb Z[x]/(x^2+x+1).
	$$
	\begin{enumerate}[(a)]
		\item Answer the following questions with suitable justification.
			\begin{enumerate}[(i)]
				\item Is $R$ a Noetherian ring?
				\item Is $R$ an Artinian ring?
			\end{enumerate}
		\item Prove that $R$ is an integrally closed domain.
	\end{enumerate}
\end{examproblem}

\begin{examproblem}[2011, 5]
	Let $R$ be a commutative ring. Recall that an element $r$ of $R$ is
	\textit{nilpotent} if $r^n=0$ for some positive integer $n$ and that
	the \textit{nilradical} of $R$ is the set $N(R)$ of nilpotent
	elements.
	\begin{enumerate}[(a)]
		\item Prove that
			$$
			N(R)=\cap_{P\text{ prime}}P.
			$$
			(Hint: given a non-nilpotent element $r$ of $R$, you may
			wish to construct a prime ideal that does not
			contain $r$ or its powers.)
		\item Given a positive integer $m$, determine the nilradical
			of $\mathbb Z/(m)$.
		\item Determine the nilradical of $\mathbb C[x,y]/(y^2-x^3)$.
		\item Let $p(x,y)$ be a polynomial in $\mathbb C[x,y]$ such that
			for any complex number $a$, $p(a,a^{3/2})=0$. Prove that
			$p(x,y)$ is divisible by $y^2-x^3$.
	\end{enumerate}
\end{examproblem}

\begin{examproblem}[2011, 6]
	Given a finite group $G$, recall that its \textit{regular representation}
	is the representation on the complex group algebra $\mathbb C[G]$ induced by
	left multiplication of $G$ on itself and its \textit{adjoint representation}
	is the representation on the complex group algebra $\mathbb C[G]$ induced
	by conjugation of $G$ on itself.
	\begin{enumerate}[(a)]
		\item Let $G=\GL_2(\mathbb F_2)$. Describe the number and dimensions of
			the irreducible representations of $G$. Then describe the
			decomposition of its regular representation as a direct
			sum of irreducible representations.
		\item Let $G$ be a group of order 12. Show that its adjoint
			representation is reducible; that is, there is
			an $H$-invariant subspace of $\mathbb C[H]$ besides 0 and
			$\mathbb C[H]$.
	\end{enumerate}
\end{examproblem}

\begin{examproblem}[2011, 8]
	Let $R$ be a commutative integral domain. Show that the following
	are equivalent:
	\begin{enumerate}[(a)]
		\item $R$ is a field;
		\item $R$ is a semi-simple ring;
		\item Any $R$-module is projective.
	\end{enumerate}
\end{examproblem}

\newpage

\section{2010}

\begin{examproblem}[2010, 1.]
	Let $p$ be a positive prime number, $\mathbb F_p$ the field with $p$ elements,
	and let $G=\text{GL}_2(\mathbb F_p)$.
	\begin{enumerate}[(a)]
		\item Compute the order of $G$, $|G|$.
		\item Write down an explicit isomorphism from $\mathbb Z/p\mathbb Z$ to
			$$
			U=\left\{
			\begin{pmatrix}
				1 & a\\
				0 & 1
			\end{pmatrix}
			\bigg|a\in\mathbb F_p\right\}.
			$$
		\item How many subgroups of order $p$ does $G$ have? Hint:
			compute $gug\inv$ for $g\in G$ and $u\in U$; use this
			to find the size of the normalizer of $U$ in $G$.
	\end{enumerate}
\end{examproblem}


\begin{examproblem}[2010, 2]
	\begin{enumerate}[(a)]
		\item Give definitions of the following terms:
			(i) a finite length (left) module, (ii) a
			composition series for a module, and (iii) the length
			of a module,
		\item Let $l(M)$ denote the length of a module $M$. Prove that if
			$$
			0\rightarrow M_1\rightarrow M_2\rightarrow\dots\rightarrow
			M_n\rightarrow 0
			$$
			is an exact sequence of modules of finite length, then
			$$
			\sum_{i=1}^n(-1)^kl(M_i)=0.
			$$
	\end{enumerate}
\end{examproblem}

\begin{examproblem}[2010, 3]
	Let $\mathbb F$ be a field of characteristic $p$, and $G$ a group of
	order $p^n$. Let $R=\mathbb F[G]$ be the group ring (group algebra)
	of $G$ over $\mathbb F$, and let $u:=\sum_{x\in G}x$ (so $u$ is an element of
	$R$).
	\begin{enumerate}[(a)]
		\item Prove that $u$ lies in the center of $R$.
		\item Verify that $Ru$ is a 2-sided ideal of $R$.
		\item Show there exists a positive integer $k$ such that $u^k=0$.
			Conclude that for such a $k$, $(Ru)^k=0$.
		\item Show that $R$ is \textbf{not} a semi-simple ring.
			(\textbf{Warning:} Please use the definition of a semi-simple
			ring: do \textbf{not} use the result that a finite length
			ring fails to be semisimple if and only if it has a non-zero
			nilpotent ideal.)
	\end{enumerate}
\end{examproblem}


\begin{examproblem}[2010, 4]
	Let $f(x)=a_nx^n+a_{n-1}x^{n-1}+\dots+a_0\in\mathbb Z[x]$
	(where $a_n\neq 0$) and let $R=\mathbb Z[x]/(f)$.
	Prove that $R$ is a finitely generated module over $\mathbb Z$ if and only if
	$a_n=\pm 1$.
\end{examproblem}

\begin{examproblem}[2010, 5]
	Consider the ring
	$$
	S=C[0,1]=\{f:[0,1]\rightarrow\mathbb R:f\text{ is continuous}\}
	$$
	with the usual operations of addition and multiplication of functions.
	\begin{enumerate}[(a)]
		\item What are the invertible elements of $S$?
		\item For $a\in[0,1]$, define $I_a=\{f\in S:f(a)=0\}$.
			Show that $I_a$ is a maximal ideal of $S$.
		\item Show that the elements of any proper ideal of $S$ have a
			common zero, i.e., if $I$ is a proper ideal of $S$, then
			there exists $a\in[0,1]$ such that $f(a)=0$ for all $f\in I$.
			Conclude that every maximal ideal of $S$ is of the form
			$I_a$ for some $a\in[0,1]$. \textbf{Hint:} as $[0,1]$
			is compact, every open cover of $[0,1]$ contains a finite
			subcover.
	\end{enumerate}
\end{examproblem}

\begin{examproblem}[2010, 7]
	Let $F$ be a field of characteristic zero, and let $K$
	be an \textit{algebraic} extension of $F$ that possesses the following
	property: every polynomial $f\in F[x]$ has a root in $K$.
	Show that $K$ is algebraically closed.\\
	\textbf{Hint:} if $K(\theta)/K$ is algebraic, consider $F(\theta)/F$
	and its normal closure; primitive elements might be of help.
\end{examproblem}


\begin{examproblem}[2010, 8]
	Let $G$ be the unique non-abelian group of order 21.
	\begin{enumerate}[(a)]
		\item Describe all 1-dimensional complex representations of $G$.
		\item How many (non-isomorphic) irreducible complex representations
			does $G$ have and what are their dimensions?
		\item Determine the character table of $G$.
	\end{enumerate}
\end{examproblem}

\newpage
\section{2009}
\begin{examproblem}[2009, 1.]
	\begin{enumerate}[(a)]
		\item Classify all groups of order $2009=7^2\times 41$.
		\item Suppose that $G$ is a group of order 2009. How many intermediate groups are
	there—that is, how many groups H are there with $1\subsetneq H\subsetneq G$, where both
	inclusions are proper? (There may be several cases to consider.)
	\end{enumerate}
\end{examproblem}

\begin{examproblem}[2009, 2]
	Let $K$ be a field. A discrete valuation on $K$ is a function $\nu:
	K\setminus\{0\}\rightarrow\mathbb Z$ such that
	\begin{enumerate}[(i)]
		\item $\nu(ab)=\nu(a)+\nu(b)$
		\item $\nu$ is surjective
		\item $\nu(a+b)\geq\text{min}\{(\nu(a),\nu(b)\}$ for
			$a,b\in K\setminus\{0\}$ with $a+b\neq 0$.
	\end{enumerate}
	Let $R:=\{x\in K\setminus\{0\}:\nu(x)\geq0\}\cup\{0\}$. Then
	$R$ is called the valuation ring of $\nu$.
	Prove the following:
	\begin{enumerate}[(a)]
		\item $R$ is a subring of $K$ containing the 1 in $K$.
		\item for all $x\in K\setminus\{0\}$, either $x$ or
			$x\inv$ is in $R$.
		\item $x$ is a unit of $R$ if and only if $\nu(x)=0$.
		\item Let $p$ be a prime number, $K=\mathbb Q$,
			and $\nu_p:\mathbb Q\setminus\{0\}\rightarrow\mathbb Z$
			be the function defined by $\nu_p(\frac ab)=n$ where
			$\frac ab=p^n\frac cd$ and $p$ does not divide $c$ and $d$.
			Prove that the corresponding valuation ring $R$ is the ring
			of all rational numbers whose denominators are relatively
			prime to $p$.
	\end{enumerate}
\end{examproblem}

\begin{examproblem}[2009, 3]
	Let $F$ be a field of characteristic not equal to 2.
	\begin{enumerate}[(a)]
		\item Prove that any extension $K$ of $F$ of degree 2
			is of the form $F(\sqrt D)$ where $D\in F$
			is not a square in $F$ and, conversely, that each such
			extension has degree 2 over $F$.
		\item Let $D_1,D_2\in F$ neither of which is a square in $F$.
			Prove that $[F(\sqrt{D_1},\sqrt{D_2}):F]=4$ if
			$D_1D_2$ is not a square in $F$ and is of degree
			2 otherwise.
	\end{enumerate}
\end{examproblem}

\begin{examproblem}[2009, 4]
	Let $F$ be a field and $p(x)\in F[x]$ an irreducible polynomial.
	\begin{enumerate}[(a)]
		\item Prove that there exists a field extension $K$ of $F$
			in which $p(x)$ has a root.
		\item Determine the dimension of $K$ as a vector space over $F$
			and exhibit a vector space basis for $K$.
		\item If $\theta\in K$ denotes a root of $p(x)$, express
			$\theta\inv$ in terms of the basis found in part (b).
		\item Suppose $p(x)=x^3+9x+6$. Show $p(x)$ is irreducible over
			$\mathbb Q$. If $\theta$ is a root of $p(x)$, compute
			the inverse of $(1+\theta)$ in $\mathbb Q(\theta)$.
			\end{enumerate}
\end{examproblem}

\begin{examproblem}[2009, 6]
	Fix a ring $R$, an $R$-module $M$, and an $R$-module homomorphism
	$f:M\rightarrow M$.
	\begin{enumerate}[(a)]
		\item If $M$ satisfies the descending chain condition on submodules,
			show that if $f$ is injective, then $f$ is surjective.
			(Hint: note that if $f$ is injective, so are
			$f\circ f$, $f\circ f\circ f$, etc.)
		\item Give an example of a ring $R$, an $R$-module $M$,
			and an injective $R$-module homomorphism $f:M\rightarrow M$
			which is not surjective.
		\item If $M$ satisfies the ascending chain condition on submodules,
			show that if $f$ is surjective, then $f$ is injective.
		\item Give an exampe of a ring $R$, and $R$-module $M$, and
			a surjective $R$-module homomorphism $f:M\rightarrow M$
			which is not injective.
	\end{enumerate}
\end{examproblem}

\begin{examproblem}[2009, 7]
	Let $G$ be a finite group, $k$ an algebraically closed field, and $V$
	an irreducible $k$-linear representation of $G$.
	\begin{enumerate}[(a)]
		\item Show that $\Hom_{kG}(V,V)$ is a division algebra with $k$
			in its center.
		\item Show that $V$ is finite-dimensional over $k$, and conclude
			that $\Hom_{kG}(V,V)$ is also finite dimensional.
		\item Show the inclusion $k\hookrightarrow\Hom_{kG}(V,V)$ found
			in (a) is an isomorphism. (For $f\in\Hom_{kG}(V,V)$, view
			$f$ as a linear transformation and consider $f-\alpha I$,
			where $\alpha$ is an eigenvalue of $f$).
	\end{enumerate}
\end{examproblem}
\newpage
\section{2008}
\begin{examproblem}[2008, 1]
	Let $f(x)$ be an irreducible polynomial of degree 5 over the field
	$\mathbb Q$ of rational numbers with exactly 3 real roots.
	\begin{enumerate}[(a)]
		\item Show that $f(x)$ is not solvable by radicals.
		\item Let $E$ be the splitting field of $f$ over $\mathbb Q$.
			Construct a Galois extension $K$ of degree 2 over
			$\mathbb Q$ lying in $E$ such that \textit{no} field $F$
			strictly between $K$ and $E$ is Galois over $\mathbb Q$.
	\end{enumerate}
\end{examproblem}

\begin{examproblem}[2008, 2]
	Let $F$ be a finite field. Show for any positive integer $n$ that
	there are irreducible polynomials of degree $n$ in $F[x]$.
\end{examproblem}

\begin{examproblem}[2008, 3.]
	Show that the order of the group $\text{GL}_n(\mathbb F_q)$ of invertible
	$n\times n$ matrices over the field $\mathbb F_q$ of $q$ elements is given by
	$(q^n-1)(q^n-q)\dots(q^n-q^{n-1})$.
\end{examproblem}

\begin{examproblem}[2008, 5]
	\begin{enumerate}[(a)]
		\item Let $R$ be a commutative principal ideal domain. Show that
			any $R$-module $M$ generated by two elements takes the form
			$R/(a)\oplus R/(b)$ for some $a,b\in R$. What more
			can you say about $a$ and $b$?
		\item Give a necessary and sufficient condition for two direct sums
			as in part (a) to be isomorphic as $R$-modules.
	\end{enumerate}
\end{examproblem}


\begin{examproblem}[2008, 6.]
	Let $G$ be the subgroup of $\text{GL}_3(\mathbb C)$ generated by the three
	matrices
	$$
	A=
	\begin{pmatrix}
		0 & 0 & 1\\
		0 & 1 & 0\\
		1 & 0 & 0
	\end{pmatrix},
	\quad B=
	\begin{pmatrix}
		0 & 0 & 1\\
		1 & 0 & 0\\
		0 & 1 & 0
	\end{pmatrix},
	\quad C=
	\begin{pmatrix}
		i & 0 & 0\\
		0 & 1 & 0\\
		0 & 0 & 1
	\end{pmatrix}
	$$
	where $i^2=-1$. Here $\mathbb C$ denotes the complex field.
	\begin{enumerate}[(a)]
		\item Compute the order of $G$.
		\item Find a matrix in $G$ of largest possible order (as an element
			of $G$) and compute this order.
		\item Compute the number of elements in $G$ with this largest
			order.
	\end{enumerate}
\end{examproblem}

\begin{examproblem}[2008, 7]
	\begin{enumerate}[(a)]
		\item Let $G$ be a group of (finite) order $n$. Show that any
			irreducible left module over the group algebra $\mathbb CG$
			has complex dimension at least $\sqrt n$.
		\item Give an example of a group $G$ of order $n\geq5$ and an
			irreducible left module over $\mathbb CG$ of complex dimension
			$\lfloor\sqrt n\rfloor$, the greatest integer to
			$\sqrt n$.
	\end{enumerate}
\end{examproblem}


\begin{examproblem}[2008, 8]
	Use the rational canonical form to show that any square matrix $M$ over
	a field $k$ is similar to its transpose $M^t$, recalling that $p(M)=0$
	for some $p\in k[t]$ if and only if $p(M^t)=0$.
\end{examproblem}

\newpage
\section{2007}


\begin{examproblem}[2007, 1]
	Let $K$ be a field of characteristic zero and $L$ a Galois extension
	of $K$. Let $f$ be an irreducible polynomial in $K[x]$ of degree 7 and
	suppose $f$ has no zeroes in $L$. Show that $f$ is irreducible
	in $L[x]$.
\end{examproblem}

\begin{examproblem}[2007, 2]
	Let $K$ be a field of characteristic zero and $f\in K[x]$ an
	irreducible polynomial of degree $n$. Let $L$ be a splitting field for
	$f$. Let $G$ be the group of automorphisms of $L$ which act trivially
	on $K$.
	\begin{enumerate}[(a)]
		\item Show that $G$ embeds in the symmetric group $S_n$.
		\item For each $n$, give an example of a field $K$ and polynomial
			$f$ such that $G=S_n$.
		\item What are the possible groups $G$ when $n=3$. Justify
			your answer.
	\end{enumerate}
\end{examproblem}

\begin{examproblem}[2007, 3.]
	Show there are exactly two groups of order 21 up to isomorphism.
\end{examproblem}

\newpage
\section{2006}


\begin{examproblem}[2006, 2]
	Let $K$ be the field $\mathbb Q(z)$ of rational functions in a variable $z$
	with coeffiecients in the rational field $\mathbb Q$.
	Let $n$ be a positive integer. Consider the polynomial $x^n-z\in K[x]$.
	\begin{enumerate}[(a)]
		\item Show that the polynomial $x^n-z$ is irreducible over $K$.
		\item Describe the splitting field of $x^n-z$ over $K$.
		\item Determine the Galois group of the splitting field
			of $x^5-z$ over the field $K$.
	\end{enumerate}
\end{examproblem}

\begin{examproblem}[2006, 3.]
	\begin{enumerate}[(a)]
		\item Let $p<q<r$ be prime integers. Show that a group of order
			$pqr$ cannot be simple.
		\item Consider groups of orders $2^2\cdot 3\cdot p$ where $p$
			has the values 5, 7, and 11. For each of those values
			of $p$, either display a simple group of order
			$2^2\cdot 3\cdot p$, or show that there cannot
			be a simple group of that order.
	\end{enumerate}
\end{examproblem}


\begin{examproblem}[2006, 4]
	Let $K/F$ be a finite Galois extension and let $n=[K:F]$. There is a theorem
	(often referred to as the "normal basis theorem") which states that there
	exists an irreducible polynomial $f(x)\in F[x]$ whose
	roots form a basis for $K$ as a vector space over $F$. You may assume
	that theorem in this problem.
	\begin{enumerate}[(a)]
		\item Let $G=\Gal(K/F)$. The action of $G$ on $K$ makes $K$ into
			a finite-dimensional representation space for $G$ over $F$.
			Prove that $K$ is isomorphic to the regular representation
			for $G$ over $F$.\\
			(The regular representation is defined by letting $G$ act
			on the group algebra $F[G]$ by multiplication on
			the left.)
		\item Suppose that the Galois group $G$ is cyclic and that $F$
			contains a primitive $n^{\text{th}}$ root of unity. Show that
			there exists an injective homomorphism $\chi:G\rightarrow
			F^{\times}$.
		\item Show that $K$ contains a non-zero element $a$ with the
			following property:
			$$
			g(a)=\chi(g)\cdot a
			$$
			for all $g\in G$.
		\item If $a$ has the property stated in (c), show that $K=F(a)$ and
			that $a^n\in F^{\times}$.
	\end{enumerate}
\end{examproblem}


\begin{examproblem}[2006, 5.]
	Let $G$ be the group of matrices of the form
	$$
	\begin{pmatrix}
		1 & a & b\\
		0 & 1 & c\\
		0 & 0 & 1
	\end{pmatrix}
	$$
	with entries in the finite field $\mathbb F_p$ of $p$ element, where
	$p$ is a prime.
	\begin{enumerate}[(a)]
		\item Prove that $G$ is non-abelian.
		\item Suppose $p$ is odd. Prove that $g^p=I_3$ for all $g\in G$.
		\item Suppose that $p=2$. It is known that there are exactly
			two non-abelian groups of order 8, up to isomorphism:
			the dihedral group $D_8$ and the quaternionic group.
			Assuming this fact without proof, determine which
			of these groups $G$ is isomorphic to.
	\end{enumerate}
\end{examproblem}


\begin{examproblem}[2006, 7.]
	There are five nonisomorphic groups of order 8. For
	each of those groups $G$, find the smallest positive integer n such that there is an
	injective homomorphism $\varphi: G\rightarrow S_n$.
\end{examproblem}

\newpage
\section{2005}

\begin{examproblem}[2005, 1.]
	For any group $G$ we define $\Omega(G)$ to be the image of the group
	homomorphism $\rho:G\rightarrow\Aut(G)$ where $\rho$ maps
	$g\in G$ to the conjugation automorphism $x\mapsto gxg\inv$.
	Starting with a group $G_0$, we define $G_1=\Omega(G_0)$ and
	$G_{i+1}=\Omega(G_i)$ for all $i\geq 0$. If $G_0$ is of order
	$p^e$ for a prime $p$ and integer $e\geq 2$, prove that $G_{e-1}$
	is the trivial group.
\end{examproblem}

\begin{examproblem}[2005, 2.]
	Let $\mathbb F_2$ be the field with two elements.
	\begin{enumerate}[(a)]
		\item What is the order of $\text{GL}_3(\mathbb F_2)$?
		\item Use the fact that $\text{GL}_3(\mathbb F_2)$ is a simple group
			(which you should not prove) to find the number
			of elements of order 7 in $\text{GL}_3(\mathbb F_2)$.
	\end{enumerate}
\end{examproblem}


\begin{examproblem}[2005, 3]
	Let $G$ be a finite abelian group. Let $f:\mathbb Z^m\rightarrow G$ be a
	surjection of abelian groups. We may think of $f$ as a homomorphism
	of $\mathbb Z$-modules. Let $K$ be the kernel of $f$.
	\begin{enumerate}[(a)]
		\item Prove that $K$ is isomorphic to $\mathbb Z^m$.
		\item We can therefore write the inclusion map
			$K\rightarrow\mathbb Z^m$ as $\mathbb Z^m\rightarrow\mathbb Z^m$
			and represent it by an $m\times m$ integer matrix
			$A$. Prove that $|\det A|=|G|$.
	\end{enumerate}
\end{examproblem}

\begin{examproblem}[2005, 4]
	Let $R=C([0,1])$ be the ring of all continuous real-valued functions
	on the closed interval $[0,1]$, and for each $c\in[0,1]$, denote by
	$M_c$ the set of all functions $f\in R$ such that $f(c)=0$.
	\begin{enumerate}[(a)]
		\item Prove that $g\in R$ is a unit if and only if $g(c)\neq0$
			for all $c\in[0,1]$.
		\item Prove that for each $c\in[0,1]$, $M_c$ is a maximal ideal of
			$R$.
		\item Prove that if $M$ is a maximal ideal of $T$, then
			$M=M_c$ for some $c\in[0,1]$. (Hint:
			compactness of $[0,1]$ may be relevant.)
	\end{enumerate}
\end{examproblem}

\begin{examproblem}[2005, 5]
	Let $R$ and $S$ be commutative rings, and $f:R\rightarrow S$
	a ring homomorphism.
	\begin{enumerate}[(a)]
		\item Show that if $I$ is a prime ideal of $S$, then
			$$
			f\inv(I)=\{r\in R:f(r)\in I\}
			$$
			is a prime ideal of $R$.
		\item Let $N$ be the set of nilpotent elements of $R$:
			$$
			N=\{r\in R:r^m=0\text{ for some }m\geq 1\}.
			$$
			$N$ is called the \textit{nilradical} of $R$. Prove that
			it is an ideal which is contained in every prime ideal.
		\item Part (a) lets us define a function
			$$
			f^*:\{\text{prime ideals of }S\}\rightarrow
			\{\text{prime ideals of }R\}.
			$$
			$$
			I\mapsto f\inv(I).
			$$
			Let $N$ be the nilradical of $R$. Show that if
			$S=R/N$ and $f:R\rightarrow R/N$ is the quotient map,
			then $f^*$ is a bijection
	\end{enumerate}
\end{examproblem}





\begin{examproblem}[2005, 7]
	Consider the polynomial $f(x)=x^{10}+x^5+1\in\mathbb Q[x]$ with
	splitting field $K$ over $\mathbb Q$.
	\begin{enumerate}[(a)]
		\item Determine whether $f(x)$ is irreducible over $\mathbb Q$ and find
			$[K:\mathbb Q]$.
		\item Determine the structure of the Galois group
			$\Gal(K/\mathbb Q)$.
	\end{enumerate}
\end{examproblem}

\begin{examproblem}[2005, 8]
	For each prime number $p$ and each positive integer $n$, how many
	elements $\alpha$ are there in $\mathbb F_{p^n}$ such that
	$F_p(\alpha)=F_{p^6}$?
\end{examproblem}

\begin{comment}
\section{Other Problems}
The exercises in Dummit and Foote section 4.4, 4.5, and 5.5 seem to cover
everything that could come up related to classifying groups of a given order.
In particular, most of the exercises in 5.5 seem quite important. Also, exercises
12-41 in Lang.

\begin{examproblem}[Dummit and Foote 5.5.6.]
	Assume that $K$ is a cyclic group, $H$ is an arbitrary group, and $\varphi_1$
	and $\varphi_2$ are homomorphisms from $K$ into $\Aut(H)$ such that
	$\varphi_1(K)$ and $\varphi_2(K)$ are conjugate subgroups
	of $\Aut(H)$. Prove by constructing an explicit isomorphism that
	$H\rtimes_{\varphi_1}K\cong H\rtimes_{\varphi_2} K$.
	[Suppose $\sigma_{\varphi_1}(K)\sigma\inv=\varphi_2(K)$
	so that for some $a\in\mathbb Z$ we have $\sigma\varphi_1(k)\sigma\inv
	=\varphi_2(k)^a$ for all $k\in K$. Show that the map $\psi:H
	\rtimes_{\varphi_1}K\rightarrow H\rtimes_{\varphi_2}K$
	defined by $\psi((h,k))=(\sigma(h),k^a)$ is a homomorphism.
	Show $\psi$ is bijective by construcing a 2-sided inverse.]
\end{examproblem}
\end{comment}

\end{document}
